\documentclass{article}
\usepackage{amsmath,amsthm}
\newcommand{\R}{\mathbb{R}}
\DeclareMathOperator{\tr}{tr}
\newtheorem{theorem}{Theorem}[section]
\title{A Short Note on Quadratic Forms}
\author{A. Author}

\begin{document}
\maketitle

\begin{abstract}
We study quadratic forms $q(x) = x^\top A x$ on $\R^n$.
\end{abstract}

\section{Introduction}\label{sec:intro}
Let $A \in \R^{n \times n}$ be \emph{symmetric}. % a comment
The trace $\tr A$ is \textbf{invariant}\footnote{Under similarity.}; see
Section~\ref{sec:main} and~\cite{book}.

\begin{equation}\label{eq:energy}
  E(x) = \frac{1}{2} x^\top A x - b^\top x
\end{equation}

\subsection{Main result}\label{sec:main}
\begin{itemize}
  \item First point, with $\lambda_1 \le \dots \le \lambda_n$.
  \item Second point, using \texttt{numpy}.
\end{itemize}

\begin{theorem}[Rayleigh]\label{thm:rayleigh}
For all $x \ne 0$, $\lambda_1 \le \frac{x^\top A x}{x^\top x} \le \lambda_n$; compare \eqref{eq:energy}.
\end{theorem}

\begin{proof}
Diagonalize $A$.
\end{proof}

\begin{align}
  a &= b + c \notag \label{eq:pending} \\
  d &= e
\end{align}

\begin{verbatim}
x = { $ % literal
\end{verbatim}

\begin{tabular}{lr}
Name & Value \\ \hline
alpha & 1 \\
\end{tabular}

Unknown \foo{kept} and Theorem~\ref{thm:rayleigh}, equation~\eqref{eq:pending}.

\begin{thebibliography}{9}
\bibitem{book} A. Writer. \emph{A Book}.
\end{thebibliography}
\end{document}
